\documentclass[11pt,a4paper]{article}
\usepackage[latin1]{inputenc}
\usepackage{amsmath}
\usepackage{amsfonts}
\usepackage{amssymb}
\usepackage{graphicx}
\author{Evac Planning Group}
\title{Notes Document}
\begin{document}
	\maketitle
	
	
	\section{Description}
	This is a living document which holds notes for the Evacuation Planning project - research conducted under the watchful eyes of Dr. Robert Heckendorn, at the University of Idaho.

	\section{Simulation} 
	\subsection{Concept} The simulation's goal is to model traffic moving around a defined city. The city is defined using the SLANG specification created by Dr. Heckendorn and Damian Ball. 
	
	The implementation uses Flex and Bison to specify a lexer and parser for SLANG. Back-end code is then implemented in C/C++. 
	
	\subsection{Work Completed}
	Currently, the work completed includes a preliminary simulation. The simulation, written in C++, uses a priority queue and a link flow formula to model traffic moving around the given city/graph.
	
	In short, the simulation computes routes for agents in the priority queue. The agents next to move, based on time, choose their next route, probabilistically, from the node they have arrived at. The amount of time it should take the agent to arrive at the next location is computed using the link flow formula, and the agent is placed into the priority queue at the appropriate time. This is accomplished using the C++ STL priority queue, which sorts the queue using a heap.
	
	The simulation runs as part of a generational genetic algorithm now. 
	
	Added safety to SLANG and necessary code to support the change to SLANG in our libraries.
	
	Added the ability to print so-called snapshots of the city in SLANG. The method is a part of the City class, and will require additional work as our project grows.
	
	
	
	\section{Misc.}
	\subsection{SLANG implementation}
	Newlines have been added to the SLANG implementation, instead of some other terminator. This means that any test files need to have an extra line at the end of the file, or risk undefined behavior.
	
	\emph{safety} has been added to the slang implementation in several places. First - edges are constructed in the grammar with a new FLOAT field representing safety. Second, a safety command has been added to specify nodes' safety. Further, a command to set safety of an edge has been added as well. Now SLANG can specify the safety of a specified node OR edge, as desired, and not only at the creation of the object.
	
	
	\section{TODO}
	Safety mutator functions should be added to the node/edge classes to make the safety of those nodes more accessible.
	
	SLANG should be modified, regarding the safety keyword - the specification should say something along the lines of 
	\begin{verbatim}
	safety: SAFETY NODE ID FLOAT NL | SAFETY EDGE ID ID FLOAT NL;
	\end{verbatim}
	where the first rule specifies the safety for a node with the given ID, and the second rule specifies the safety of the edge from the node with the first ID to the node with the second ID. 
	
	We need some sort of visualization. Currently considering the R programming language to achieve visualization, however, that will wait a while longer.
	
	A new, more in-depth test file is also needed, however small tests are sufficient to some extent currently.
\end{document}